% ATTEMPT_STATUS: unresolved
% RESEARCH_GENERATED: 2026-10-01; GPT-6 Astra Ultra; individual mathematical attempt
% TOP_PROBLEM: {"id": 159, "title": "Uniform Random Variables with Uniform Sum", "category_id": 2, "category": {"id": 2, "name": "combinatorics", "display_name": "Combinatorics", "description": "Counting problems, graph theory, discrete structures.", "slug": "combinatorics", "order_index": 2, "created_at": "2026-07-31T15:26:25.671Z"}, "status": "open"}
% FIRST_POSTED: 2026-10-01
% Imported from: unsolvedmath_additional_500.tex; UnsolvedMath ID 159
% !TeX program = lualatex
% Compile twice: lualatex 159.tex
% Self-contained: no external bibliography, images, or \input files.
% Numeric section labels and visible numbers are original UnsolvedMath IDs.
\documentclass[11pt,a4paper]{article}
\usepackage[margin=24mm,headheight=14pt]{geometry}
\usepackage{fontspec}
\setmainfont{Latin Modern Roman}
\setsansfont{Latin Modern Sans}
\setmonofont{Latin Modern Mono}
\usepackage{amsmath,amssymb,mathtools}
\usepackage{unicode-math}
\setmathfont{Latin Modern Math}
\usepackage{microtype}
\usepackage{enumitem,xurl,needspace}
\usepackage[dvipsnames]{xcolor}
\usepackage{hyperref}
\hypersetup{unicode=true,colorlinks=true,linkcolor=MidnightBlue,urlcolor=MidnightBlue,
 pdftitle={UnsolvedMath Problem 159},pdfauthor={Source-annotated compilation},bookmarksnumbered=true}
\usepackage{bookmark,tocloft,titlesec,fancyhdr}
\setlength{\cftsecnumwidth}{4.2em}
\cftsetpnumwidth{2.5em}
\cftsetrmarg{4em}
\titleformat{\section}{\Large\bfseries\raggedright}{\thesection}{0.7em}{}
\pagestyle{fancy}\fancyhf{}
\fancyhead[L]{\small\sffamily UnsolvedMath research attempt}
\fancyhead[R]{\small Original catalogue IDs}
\fancyfoot[C]{\thepage}
\setlength{\parindent}{0pt}
\setlength{\parskip}{5pt plus 1pt}
\setlength{\emergencystretch}{3em}
\setcounter{tocdepth}{1}\setcounter{secnumdepth}{1}
\setlist[itemize]{leftmargin=3em,itemsep=4pt,topsep=4pt,parsep=0pt}
\allowdisplaybreaks
\newcommand{\N}{\mathbb{N}}
\newcommand{\Z}{\mathbb{Z}}
\newcommand{\Q}{\mathbb{Q}}
\newcommand{\R}{\mathbb{R}}
\newcommand{\C}{\mathbb{C}}
\newcommand{\F}{\mathbb{F}}
\newcommand{\A}{\mathbb{A}}
\newcommand{\PP}{\mathbb{P}}
\newcommand{\E}{\mathbb{E}}
\newcommand{\eps}{\varepsilon}
\newcommand{\dd}{\,\mathrm d}
\newcommand{\sref}[2]{\hyperlink{source:#1:#2}{[S#2]}}
\newif\ifshowcatalogue
\showcataloguetrue
\newcommand{\cataloguescope}[1]{\ifshowcatalogue
 \paragraph{Statement recorded in the supplied source.}
 #1\fi}
\begin{document}
% UnsolvedMath ID 159; source catalogue code GREEN-071

\setcounter{section}{158}

\section{Uniform Convolutions of Integer-Valued Distributions}\label{159}

{\small\sffamily UnsolvedMath ID 159 · Combinatorics}

\subsection{Definitions and mathematical statement}

\paragraph{Shared notation.}Unless a section says otherwise, $\N=\{1,2,\ldots\}$,
$\Z$ is the ring of integers, $\Q,\R,\C$ are the rational, real, and complex
fields, $[N]=\{1,\ldots,\lfloor N\rfloor\}$, and $\log$ is the natural
logarithm. For real functions of a parameter tending to infinity,
$f\ll g$ and $f=O(g)$ mean $|f|\leq Cg$ eventually for some fixed $C$;
$f\gg g$ reverses this comparison; $f=o(g)$ means $f/g\to0$;
$f\sim g$ means $f/g\to1$; and $f\asymp g$ means both order bounds.
A subscript on $O$ or $\ll$ specifies allowed constant dependence.
The expression $n^{o(1)}$ in an upper bound means growth smaller than
$n^\epsilon$ for each fixed $\epsilon>0$. Local definitions govern any
exception, finite-field convention, or use of the same symbol for another
object. An asymptotic claim is not automatically an effective algorithm.



Let $p_i=\Pr(X=i)$ and $q_j=\Pr(Y=j)$ have finite supports. Independence gives $\Pr(X+Y=s)=\sum_i p_iq_{s-i}$. A distribution is uniform on its support if all its nonzero probabilities are equal. The proposed implication is that a uniform convolution forces both input distributions to be uniform on their respective supports. \sref{159}{1} \sref{159}{2}

\paragraph{Statement recorded in the supplied source.}
Suppose $X, Y$ are finitely-supported independent random variables taking integer values such that $X + Y$ is uniformly distributed on its range. Are $X$ and $Y$ themselves uniformly distributed on their ranges? \sref{159}{1}

\paragraph{Residual task reported by the archive.}

The embedded review (2026-08-17) records: Prove the unfair 0-1 polynomial conjecture for arbitrary degree or construct a nonnegative-real counterexample. Computational verification cannot settle the unbounded-degree statement. \sref{159}{1}

\subsection{Short English statement}

If the sum of two independent, finitely supported integer-valued random variables is uniform, must each original variable also be uniform?

\subsection{Research attempt}

\paragraph{Provenance and scope.}
This individual attempt was generated by GPT-6 Astra Ultra on 1 October 2026. The method was to derive explicit partial results and test the first proposed extension adversarially. The original statement and sources below are retained. No claim of mathematical novelty or complete solution is made.

\paragraph{Formal target and primary-source context.}
The variables are independent, finitely supported, and integer-valued. Uniformity is on their actual supports, which need not be intervals. The inspected work of Ghidelli and the computational follow-up treat the equivalent nonnegative polynomial-factor question and restricted cases, rather than an unrestricted proof. We give the precise normalisation, prove the target when one variable has at most two possible values, and identify the coefficient recurrence that stops the same proof in larger support.

\paragraph{The polynomial translation.}
Translate each variable so its minimum value is zero. Let $p_i=\Pr(X=i)$ and $q_j=\Pr(Y=j)$, with $p_0,q_0>0$. If $X+Y$ is uniform on its support, its mass there must be $p_0q_0$, since zero has the unique representation $0+0$. Put
\[
 P(z)=\sum_i(p_i/p_0)z^i,\qquad Q(z)=\sum_j(q_j/q_0)z^j.
\]
Both have nonnegative real coefficients and constant term one. Independence says that $P(z)Q(z)$ has every coefficient equal to zero or one. Conversely any such factorisation, after normalising both factors at $z=1$, gives independent finite variables with a uniform sum. Uniformity of each factor variable is exactly the assertion that all its nonzero normalised coefficients equal one. Thus the algebraic formulation preserves all probability normalisations.

\paragraph{A complete proof for a two-point factor.}
Suppose the support of $X$ is $\{0,a\}$ with $a>0$. Then $P(z)=1+rz^a$ for some $r>0$. In the product, the coefficient of $z^a$ is $r+q'_a$, where $q'_a\geq0$ is the coefficient of $Q$. This coefficient is positive and at most one, so $r\leq1$.

Now reverse both polynomials and divide each reversal by its constant coefficient, which is the leading coefficient of the original polynomial. Since the original product has leading coefficient one, the product of the two reversed normalised factors is again a zero-one polynomial. The reversed normalised first factor is $1+r^{-1}z^a$. The same coefficient argument therefore gives $r^{-1}\leq1$. Hence $r=1$, proving that $X$ is fair on its two-point support.

It remains to prove fairness of $Y$, rather than assume it. Write $Q(z)=\sum b_jz^j$ and $P Q=\sum e_jz^j$, with $e_j\in\{0,1\}$ and coefficients outside support interpreted as zero. Then
\[
 b_j+b_{j-a}=e_j.
\]
Along each residue class modulo $a$, start below zero with coefficient zero and proceed upward. If the previous coefficient is zero or one, then the next is $e_j-b_{j-a}$; nonnegativity forces it again to be zero or one. Induction proves that every $b_j\in\{0,1\}$. Thus $Y$ is uniform too. If $X$ has only one support point, the conclusion is immediate from translation of the sum, so all cases with support size at most two are solved.

\paragraph{A structural consequence for fair inputs.}
If both inputs are uniform on supports $A,B$ of sizes $m,k$, then the probability of a sum $s$ is $r(s)/(mk)$, where $r(s)$ counts ordered pairs from $A\times B$. The smallest sum has one representation, so a uniform output forces $r(s)=1$ for every represented $s$. Equivalently $|A+B|=mk$. Thus any positive resolution of the general conjecture would also force a direct-sum support structure. Uniformity of the factors alone does not guarantee uniformity of the output: for two fair Bernoulli variables, the output masses are one quarter, one half, one quarter.

\paragraph{Where the recurrence argument stops.}
For a larger factor $P=1+\sum_{i>0}a_iz^i$, the coefficient equation becomes $b_j=e_j-\sum_{i>0}a_ib_{j-i}$. Nonnegativity bounds the entire weighted sum, but does not separately force every $a_i$ or $b_j$ to be zero or one. The reverse-polynomial trick controls the first and last nonzero terms but leaves interior interactions. Declaring the recurrence integer-valued would be circular: its coefficients are the unknown real numbers whose integrality-like property is the question.

\paragraph{Verification and status.}
Reversal preserves nonnegativity, and its normalising constants multiply to one because the product's leading coefficient is one. The recurrence uses the actual gap $a$, so nonconsecutive supports are covered. The proof is exact and requires no numerical root search. The target is established here for every two-point factor and for degenerate variables. The general multi-point factorisation remains unresolved; the first missing step is a coefficient argument controlling interior real weights without presupposing that they are already zero or one.

\subsection{Sources}

{\small

\begin{itemize}

\item[\hypertarget{source:159:1}{S1}] ULAM.AI, \emph{UnsolvedMath}, supplied \texttt{problems.json}, record 159 (GREEN-071), statement and background. The embedded literature-review date is 2026-08-17; that is the archive’s review date, not a fresh certification. Dataset landing page: \url{https://huggingface.co/datasets/ulamai/UnsolvedMath}.

\item[\hypertarget{source:159:2}{S2}] Luca Ghidelli, Progress on the unfair 0-1-polynomials conjecture using linear recurrences and numerical analysis, arXiv:2209.09843 (2022). \url{https://arxiv.org/abs/2209.09843}. Bibliographic information imported from the supplied record.

\item[\hypertarget{source:159:3}{S3}] Kevin G. Hare, Computational Progress on the Unfair 0-1 Polynomial Conjecture, arXiv:2307.07363; Experimental Mathematics, published online 2025. \url{https://arxiv.org/abs/2307.07363}. Bibliographic information imported from the supplied record.

\item[\hypertarget{source:159:4}{S4}] David Zhang, research report answering MathOverflow question 339137, 16 June 2026. \url{https://mathoverflow.net/questions/339137/why-do-polynomials-with-coefficients-0-1-like-to-have-only-factors-with-0-1}. Bibliographic information imported from the supplied record.

\item[\hypertarget{source:159:5}{S5}] Anatoly Zhigljavsky, Nina Golyandina and Svyatoslav Gryaznov, Deconvolution of a discrete uniform distribution, Statistics \& Probability Letters 118 (2016), 37-44. \url{https://doi.org/10.1016/j.spl.2016.06.006}. Bibliographic information imported from the supplied record.

\item[\hypertarget{source:159:6}{S6}] Ben Green, 100 Open Problems, Problem 28, most recent update December 2025. \url{https://people.maths.ox.ac.uk/greenbj/papers/open-problems.pdf}. The author-hosted document was inspected on 27 September 2026; the archive’s local code is not a problem-number locator in this paper.

\item[E1] L. Ghidelli, \emph{Progress on the unfair 0-1-polynomials conjecture using linear recurrences and numerical analysis}. \url{https://arxiv.org/abs/2209.09843}. Inspected 1 October 2026.

\item[E2] \emph{Computational progress on the unfair 0-1 polynomial Conjecture}, primary manuscript. \url{https://arxiv.org/abs/2307.07363}. Inspected 1 October 2026.

\end{itemize}

}

\label{end:159}
\end{document}
